The fourth chapter changes the nature of the question. Lines and planes arose by
restricting motion to fixed directions. Conic sections arise when a moving point
is required to satisfy a relation between distances. Because Euclidean distance
contains squares and square roots, translating such requirements into Cartesian
language naturally produces equations of degree two.

The philosophical intention is to reverse the usual order of presentation. The
curve is not introduced as a formula whose graph must be memorised. It is first
created by a geometric demand. The equation is then derived as the coordinate
record of that demand. Only after the objects are understood geometrically do we
ask how to recognise them when a quadratic equation is given first.
